\documentclass[11pt]{article}
\usepackage[margin=25mm]{geometry}
\usepackage{amsmath,amssymb}
\usepackage{booktabs}
\begin{document}

\title{I01 prevalidation 12 result:\\Selection rule for paired coherent constraint leakage}
\date{2026-10-04}
\maketitle

\section*{Exact-Aer atlas}
At $\epsilon=1$ rad, the $P=0$, $\tau=0$ constraint contrast has the following
maximum pooled values:
\begin{center}
\begin{tabular}{lr}
\toprule
intervention & constraint contrast\\
\midrule
one-sided $R_x$ on qubit 0 & $5.6\times10^{-17}$\\
same-qubit inverse $R_x$ & $1.44217\times10^{-3}$\\
same-qubit inverse $R_y$ & $1.44217\times10^{-3}$\\
same-qubit inverse $R_z$ & $2.8\times10^{-17}$\\
$R_x$ inverse pair from qubit 0 to 1 & $4.9\times10^{-17}$\\
same-qubit equal-sign $R_x$ & $1.44218\times10^{-3}$\\
\bottomrule
\end{tabular}
\end{center}

\section*{Selection rule within this six-row atlas}
Within this initial atlas, leakage requires both a preparation-side and
post-science-layer operation and a sector-mixing transverse axis ($X$ or $Y$).
A diagonal $Z$ operation, one-sided operation, or the particular return operation
from qubit zero to its conjugate partner qubit one leaves the constraint null exact.
Changing the second transverse pulse from inverse to equal sign reverses the
contrast orientation while preserving its magnitude to this precision.

\section*{Scope}
This is a circuit-level mechanism map. Its original same-qubit necessity inference
was intentionally tested next and is not retained: stage 13 finds nonzero transfers
from qubit zero to other momentum pairs. It remains a toy-model result, not evidence
for physical dimension selection.

\section*{Reproduction}
Code: \texttt{../src/relational\_atlas\_i01/holdout/qiskit\_rank\_holdout12\_intervention\_atlas.py}.\\
Output: \texttt{../results/qiskit\_rank\_holdout12\_intervention\_atlas.json}.
\end{document}
